\section{Mechanism of Forking}
\label{sec:mechanism}

A finite corpus shows each context only some of its plausible continuations, so a validation position can pair a \emph{seen context} with a \emph{novel continuation} that never followed it in training. Replay reinforces the observed continuations without supplying evidence for the others, and can crowd those alternatives out. We show how this happens, then test why the $n$-gram module allows it to proceed separately within each context; additional experimental details are in Appendix~\ref{app:mechanistic}.

\begin{figure}[!t]
\centering
\begin{subfigure}[b]{0.58\linewidth}
\centering
\includegraphics[width=\linewidth]{figures/walk_their_dog_loss.pdf}
\subcaption{}
\label{fig:walk-their-dog}
\end{subfigure}\hfill
\begin{subfigure}[b]{0.38\linewidth}
\centering
\includegraphics[width=\linewidth]{figures/fig09_crowd_ce_seen_novel.pdf}
\subcaption{}
\label{fig:crowd-ce}
\end{subfigure}
\caption{Crowding-out: (a) context-only continuation losses for ``walk their dog''; (b) validation cross-entropy of train-seen and novel continuations, pooled over contexts.}
\label{fig:crowding}
\end{figure}

\subsection{Crowding-out under replay}
\label{sec:sharpening}
Forking is ordinary overfitting confined to very small problems. For a context $c_i$, call the next-token logits computed from it its \emph{context subspace}, and split the vocabulary into the continuations of $c_i$ seen in training, those that appear only in validation, and the rest (Appendix~\ref{app:crowding-derivation}). If training shows a single continuation $C$, the loss in that subspace is $-\ln p_C$; an idealized step directly in logit coordinates raises the logit of $C$ by $\eta(1-p_C)$ and, through the softmax, lowers every other logit by $\eta p_v$. Mass flows from the unseen continuations onto the seen one, so the seen-to-unseen odds grow by a factor above one at every pass with nothing to stop them while $p_C<1$ (Appendix Equation~\eqref{eq:zt}). We call this \emph{crowding-out}; it provides an idealized account of the growth in the novel-continuation penalty with replay.

Figure~\ref{fig:walk-their-dog} shows the picture on one context. ``walk their dog'' is followed by \emph{down} once in training and, in the fixed validation window, once each by \emph{daily} and \emph{for}. Probed with context-only input, the probability of \emph{down} rises by more than an order of magnitude over five passes while \emph{daily} and \emph{for} are compressed---\emph{daily} by orders of magnitude---and so is the total mass outside the continuations observed in either sample. Figure~\ref{fig:crowd-ce} shows that the context is typical. Pooling validation positions whose bigram context occurs in training, split by whether the continuation was seen there (405{,}058 positions) or not (158{,}790), both models do well on seen continuations over the first three passes: the control's cross-entropy falls from 4.3 to 2.8 nats and the both-table model's stays near 3.4. On novel continuations they diverge: the control still improves, 6.9 to 5.6 nats, while the both-table model deteriorates from 7.5 to 9.0 with a step at every epoch boundary---the staircase of Figure~\ref{fig:staircase}. Consistently, at epoch~3 the both-table model's predictive entropy on validation is 2.93 nats against a cross-entropy of 5.26, whereas the control's are 3.59 and 3.58 (same runs as Figure~\ref{fig:crowd-ce}): far more confident than accurate, the signature of a model that has committed to few continuations per context \citep{carlsson2025hyperfitting}. Crowding-out also explains why freezing the tables does not stop forking (Section~\ref{sec:interventions}): the subspace coordinate is already in place, and the backbone can continue the descent on its own.

Crowding-out hurts more where novel continuations are common. Rare contexts have had fewer chances to reveal their continuations; under heavy-tailed conditional distributions, the novel fraction falls slowly with frequency, offering an explanation for the approximate power-law trend in the gap. Appendix~\ref{app:math} develops this connection through the seen/novel gap decomposition and Good--Turing estimation, including the terms that prevent exact proportionality.

\subsection{Context subspaces are independent}
\label{sec:context-locality}
\begin{figure}[!t]
\centering
\begin{subfigure}[b]{0.45\linewidth}
\includegraphics[width=\linewidth]{figures/fig07a_context_mask.pdf}
\subcaption{Gap contribution.}
\label{fig:ctxmask-gap}
\end{subfigure}\hfill
\begin{subfigure}[b]{0.45\linewidth}
\includegraphics[width=\linewidth]{figures/fig07b_seen_mass.pdf}
\subcaption{Mass on train-seen continuations.}
\label{fig:ctxmask-mass}
\end{subfigure}
\caption{Random context loss masking.}
\label{fig:context-subspace-mask}
\end{figure}
Section~\ref{sec:sharpening} treats each context subspace as evolving on its own, with the sharpening within it as the engine of forking. Both statements are claims about what the $n$-gram module buys: a private coordinate per context, so that updates driven by one context barely touch the conditional distribution of another. We test this by ablation. We hold out a random 2\% of training trigram contexts and exclude the next-token loss at their positions from the start of epoch~2 or epoch~3 onward, keeping inputs, trainable parameters and the loss denominator unchanged; a disjoint 2\% serves as unmasked peers, and each masked run is paired with normal training from the same checkpoint. Figure~\ref{fig:ctxmask-gap} plots their gap contribution $C_S=\sum_{i\in S}q_i[\bar\ell_{\mathrm{val},i}-\bar\ell_{\mathrm{tr},i}]$ on fixed batches. Until the mask fires the masked curves track normal training; afterwards they freeze---for two and three further passes---while the unmasked and normal curves take a fresh jump at every replay boundary. The gap of a masked context stops growing the moment its own replay stops, although the remaining 98\% of contexts keep training throughout: gap growth is generated inside each subspace, not imported from the others.

Figure~\ref{fig:ctxmask-mass} follows the sharpening coordinate itself, the predictive mass on the continuations of each context that training has seen. In the n-gram model the unmasked contexts pour mass into their seen continuations pass after pass ($0.43\to0.69$ by epoch~4) while the masked ones freeze ($0.43\to0.40$): once a context's own loss is removed, nothing else sharpens its subspace. Without the $n$-gram module the mask changes nothing ($0.22\to0.39$ masked vs.\ $0.41$ unmasked), and replay raises the mass on seen and validation-only continuations alike ($0.013\to0.06$): every context's conditional distribution lives in parameters shared with all others, and repetition improves it holistically. The table therefore partitions optimization into weakly interacting context subspaces; inside a subspace, replay-driven sharpening proceeds unchecked by the continuations of other contexts, which is exactly the over-memorization assumed in Section~\ref{sec:sharpening}. Masked contexts are not immutable---the shared backbone still drifts---but that drift moves train and validation losses together and never enters the gap of Figure~\ref{fig:ctxmask-gap}.

